%% =====================================================================================
%% 8 Discussion, threats to validity, limitations, ethics.
%% Do not promise reproducibility of model text: replaying the same cell with the same
%% seed gives different play. Promise the design (prompts, configurations, seeds, code).
%% The "must"/unstated-goal threat is no longer untested: the neutral-wording and
%% goal-stated arms answer it (Section 4). Keep the remaining threats honest.
%% =====================================================================================
\section{Discussion}
\label{sec:discussion}

No model grades its payments by the size of the risk under the base prompt. Four models pay
a default of their own, close to the fair share or well above it, and still pay at $p=0$,
where paying helps no one, so generosity alone cannot explain their play. GPT Luna stops at
$p=0$ but pays the same share at every positive risk. The default does not come from
normative wording. It appears when the goal is left unstated, and stating the goal removes
most payment at zero risk without making most models weigh the risk. All five adjust to
partners only in simple ways, mostly by stopping once the outcome is settled.

\paragraph{Possible causes.}
Our data record choices, not reasoning, so we offer hypotheses. First, the equal split may be
a focal point, a choice everyone expects everyone else to make; models rebuild it without the
hint, and a 50/50 split can likewise reflect a focal point rather than a
preference~\cite{dodivers2025uncovering}. Second, a question may call up computation and a
move request a default, which fits reports that models can hold a correct belief or plan and
act against it~\cite{fan2024can,schmied2026llms}. Third, post-training may favour doing one's
share when nothing says otherwise, which fits the finding that models are more cooperative
than people~\cite{mei2024turing} and our finding that naming the goal removes most payment at
zero risk. Stronger reasoning is not an obvious fix, since reasoning models free-ride more in
public goods games~\cite{guzmanpiedrahita2025corrupted}.

\paragraph{Consequences for design.}
\emph{State the agent's goal.} An agent that is told only the rules falls back on a default,
and one sentence naming its goal removed the payment at zero risk that neutral wording left in
place. Evaluations should report what the agent was told to want.
\emph{Test against the optimum, not the success rate.} A group can reach its target in every
game and still waste most of its money, so a study needs $\mathrm{opt}(p)$, a condition where
paying is dominated, such as $p=0$, and a check for payment after the outcome is settled.
\emph{Test in mixed populations.} Self-play hides how defaults interact: an exact fair-share
default is optimal among copies of itself and fragile next to one late dropout, while a
generous default absorbs that partner and pays for others. \emph{Choose the selection signal
with care.} Any leaderboard, market or imitation process that scores an agent against its
tablemates favours whoever pays least, by Proposition~\ref{prop:table}, and in our games this
lowered group success, while scoring across independent groups favoured the agent whose
groups succeed. Each test needs only the game rules, a few scripted policies or a second
model, and each caught a failure that self-play success rates did not show.

\paragraph{Threats to validity.}
The goal sentence is itself an instruction, and we test one phrasing of it, at three risk
levels, in self-play only. Two other features of the setup may push models toward paying, and
neither is tested here. The decision prompt asks for a final line with the move, so models
may deliberate less when deciding than when answering a question. And the game is well
known, so the fair share may be recalled rather than derived; the neutral wording removes the
game's name, but the rules still describe it. The same prompts were used for every model, so
these features cannot explain why the models differ, but they may raise the common level of
payment.

\paragraph{Limitations.}
We test five low-cost models, one per provider, in English and in one game; more capable
models may respond to risk~\cite{anon-if}, so our claims are about this tier. Cells hold ten
games, which limits power for the noisiest model, Grok. Our rewordings test wording, not
every prompt design. The welfare losses assume risk-neutral players; the two tests where
paying is dominated do not. The scripted partners, mixed tables and selection analysis use the
base prompt only; mixed tables hold two models with fixed seat blocks within a cell, and
$\alpha$-Rank assumes rare entry. For four models the self-play grid was collected separately;
a rerun of the same prompt matches it (Table~\ref{tab:robust}). Model text is not
reproducible, so we release prompts, configurations, seeds and code, and promise
reproducibility of the design rather than of the transcripts.

\paragraph{Ethics.}
The study involves no human subjects and no personal data. Its findings describe failure
modes of deployed systems, and we report them so that designers can test for them before
agents are trusted with shared resources.
